% =====================================================================
%  c01 —— §1 引言（Lead 独占）
%  对应原文 PDF pp.2--3（印刷页码 2--3，偏移 0）
%  本节无编号陈述；p.3 的两条 "Theorem." 在原文中即为**不编号**的定理段，
%  用 thmplain 排版，切勿误编为 Theorem 1/2（编号 1 从 p.9 才开始）。
% =====================================================================

% [原文 p.2]
\section{引言}

\para{论文。}
\begin{itemize}[leftmargin=1.6em,itemsep=2pt,topsep=2pt]
\item \emph{Simple geodesics and Weil--Petersson volumes of moduli spaces of bordered
  Riemann surfaces.} Invent. Math., 167(1):179--222, 2007.
\item \emph{Weil--Petersson volumes and intersection theory on the moduli space of
  curves.} J. Amer. Math. Soc., 20(1):1--23, 2007.
\end{itemize}

% [Tutor Note: 这两篇（连同 [Mir08]）是 Mirzakhani 2014 年获 Fields 奖的三篇代表作。
%  第一篇解决"体积递推"，第二篇解决"系数是相交数"，第三篇（简单闭测地线计数）
%  把同一套体积递推用到计数问题上。下文 §3 的讲义主要按前两篇的顺序展开
%  （背景 → 长度恒等式 → 体积递推 → 辛约化与相交数 → Witten--Kontsevich）；
%  第三篇的计数结果不是 §3 的独立讲题。]

\para{论文的目标。}
\begin{enumerate}[leftmargin=1.8em,itemsep=2pt,topsep=2pt]
\item 导出 WP 模空间体积多项式的一条显式递推。
\item 应用辛约化，证明这些多项式的系数是相交数。
\item 证明该体积递推满足 Virasoro 关系——即 Witten 猜想。
\end{enumerate}

Maryam Mirzakhani 的这三篇论文 [Mir07a, Mir07b, Mir08] 之所以引人注目，
在于它们牵涉了大量高度发展的考量与内容。第一篇工作需要详细描述 Teichm\"uller 空间、
映射类群的作用、Weil--Petersson（WP）辛几何的公式、裤子（pair of pants）上简单
测地弧的分类、长度无穷和恒等式以及积分的精确计算。第二篇工作涉及模空间
tautological 类 $\kappa_1$ 与 $\psi$ 的描述（作为双曲几何中 $S^1$ 主丛的特征类）、
矩映射与精确的辛约化，以及组合计算。第三篇工作用到 Thurston 测度测地叶状结构
空间 $\cM\cG\cL$ 的分片线性（PL）结构、列车轨道（train-track）辛形式，以及 Masur
关于映射类群在 $\cM\cG\cL$ 上作用遍历的结果。这些工作的一个精彩之处在于，
它们展示了 Riemann／双曲曲面的几何、拓扑与形变理论的重要方面。三篇工作以及一篇
近期预印本都应用了 Mirzakhani 的体积递推。把 Mirzakhani 的工作放进 Witten--Kontsevich
理论的诸种研究途径中加以讨论，见 Mulase--Safnuk 论文 [MS08] 的第一节。Kontsevich
对 Witten 猜想的原始解法 [Kon92]（包括模空间上 tautological 类的基本几何与矩阵
积分的带状图展开）的简要阐述，见 Looijenga 的 Bourbaki 讨论班报告 [Loo93]。
Mirzakhani 体积递推、Witten--Kontsevich 的解法以及 WP 体积极限应用的简要阐述，
见 [Do11]。Mirzakhani 素简单测地线定理 [Mir08] 的整体性阐述，见 [Wlp10, 第 9、10 章]。
Mirzakhani 的工作只是这个高度活跃、研究者众多的学科的一部分；Google Scholar 显示
这三篇 Mirzakhani 论文被引用 200 次，而更一般地，Witten 关于二维引力与规范理论的
原始论文 [Wit91, Wit92] 被引用超过 1100 次。

% [Tutor Note: 原文 "Mirzakhani's prime simple geodesic theorem" 中的 "prime" 指
%  "素"（不可再分）的简单闭测地线，即在本原（primitive）意义下不可约的闭测地线，
%  不是"质数"。中译沿用"素简单测地线"这一约定。]

\para{第一篇论文简述。}
带边 Riemann／双曲曲面模空间覆盖的辛纤维积分解，与辛形式的
$\dd\ell\wedge\dd\tau$ 公式、以及测地长度和的一条万有恒等式相结合，
导出计算体积的一条显式递推。亏格 $g$、$n$ 条边界的曲面模空间的体积，
被证明是曲面边界长度的一个系数全正的多项式，其总次数为 $6g-6+2n$。

% [原文 p.3]
\para{第二篇论文简述。}
扭转 Hamilton 流与边界长度矩映射相结合，可对带边双曲曲面模空间族应用辛约化。
其推论是：体积多项式的系数是模空间的特征数。一个几何构造表明这些特征数就是
tautological 相交数。体积递推则被证明满足 Virasoro 代数约束。

下面的定理是这些论文最直接的结论。

% 原文 p.3 上两条 "Theorem." 均**不编号**
\begin{thmplain}
WP 体积多项式。体积多项式由总次数更小的体积多项式递推确定，
[Mir07a, 公式 (5.1) 与定理 8.1]。亏格 $g$、$n$ 条边界、边界长度为
$L=(L_1,\dots,L_n)$ 的双曲曲面的模空间，其体积 $V_{g,n}(L_1,\dots,L_n)$ 是多项式
\[
V_{g,n}(L)=\sum_{\substack{\alpha\\ |\alpha|\le 3g-3+n}} C_\alpha L^{2\alpha},
\]
其中 $\alpha=(\alpha_1,\dots,\alpha_n)$ 为多重指标，且 $C_\alpha>0$ 落在
$\pi^{6g-6+2n-2|\alpha|}\Q$ 中，[Mir07a, 定理 1.1 与 6.1]。其系数是相交数，由下式给出
\[
C_\alpha=\frac{2^{\delta_{1g}\delta_{1n}}}{2^{|\alpha|}\alpha!\,(3g-3+n-|\alpha|)!}
\int_{\Mgnb}\psi_1^{\alpha_1}\cdots\psi_n^{\alpha_n}\,\omega^{3g-3+n-|\alpha|},
\]
其中 $\psi_j$ 是沿第 $j$ 个穿刺点的余切线的 Chern 类，$\omega$ 是辛形式，
$\alpha!=\prod_{j=1}^n\alpha_j!$，而 $\delta_{**}$ 是 Kronecker 指示 delta，
[Mir07b, 定理 4.4]。
\end{thmplain}

% [Tutor Note: 这里 $L^{2\alpha}$ 是多重指标简写，指 $\prod_i L_i^{2\alpha_i}$。
%  分母中的 $2^{|\alpha|}\alpha!$ 不是 $\omega$ 归一化的产物：在最顶次 $|\alpha|=3g-3+n$
%  时被积式里根本没有 $\omega$（只有 $\int\psi^\alpha$），却仍带 $2^{|\alpha|}$，
%  所以它只能来自 $C_\alpha$ 与相交数 $\langle\tau^\alpha\rangle$ 之间的换算：
%  定理自身的定义给出 $(\alpha)_g=C_\alpha\,2^{-\delta_{1g}\delta_{1n}}
%  \prod_i\alpha_i!\,2^{|\alpha|}$，而在顶次
%  $(\alpha)_g=\langle\tau^{\alpha_1}\cdots\tau^{\alpha_n}\rangle_g=\int\psi^\alpha$（见 p.31），
%  故 $C_\alpha=2^{\delta_{1g}\delta_{1n}}\langle\tau^\alpha\rangle/(2^{|\alpha|}\alpha!)$。
%  （它与递归中双阶乘 $(2k+1)!!=(2k+1)!/(2^k k!)$ 的归一化一致，也可看作按矩映射变量
%  $L_j^2/2$ 展开体积的产物。）讲义取 $\omega=2\omega_{WP}=2\pi^2\kappa_1$
%  （p.4、p.6、p.10），这一条与 $2^{|\alpha|}$ 无关。
%  核对锚点：$(g,n)=(1,1)$ 时 $C_0=2\int_{\Mbar_{1,1}}\omega=2\cdot(\pi^2/12)=\pi^2/6$、
%  $C_1=\frac{2}{2}\int_{\Mbar_{1,1}}\psi=1/24$，与 $V_{1,1}=\pi^2/6+L^2/24$ 逐项吻合。]
%
% [Tutor Note: 原文写 $\delta_{**}$，这是"Kronecker 指示 delta"的缩写记法，
%  实指 $\delta_{1g}\delta_{1n}$（即只有 $(g,n)=(1,1)$ 时 $\delta_{**}=1$，其余为 0），
%  故公式在 $(1,1)$ 情形多出因子 $2$（注意不是 $1/2$）。它对应原文 p.20 的说明：
%  椭圆对合使椭圆情形的相交数乘 $1/2$（等价地，$V_{1,1}$ 按给定值的 $1/2$ 形式定义），
%  因此按"给定"体积写 $C_\alpha$ 时要在 $(1,1)$ 情形补回因子 $2$。]

\begin{thmplaincite}{Mir07a, 定理 6.3 与 6.4}
体积多项式首项系数的递推关系。对多重指标 $\alpha$，定义
\[
(\alpha_1,\dots,\alpha_n)_g=C_\alpha\times 2^{-\delta_{1g}\delta_{1n}}
\times\prod_{i=1}^n\alpha_i!\times 2^{|\alpha|},
\]
则当 $n>0$ 且 $\sum_i\alpha_i=3g-3+n$ 时，有 dilaton 方程
\[
(1,\alpha_1,\dots,\alpha_n)_g=(2g-2+n)(\alpha_1,\dots,\alpha_n)_g,
\]
而当 $n>0$ 且 $\sum_i\alpha_i=3g-2+n$ 时，有 string 方程
\[
(0,\alpha_1,\dots,\alpha_n)_g=\sum_{\alpha_i\neq 0}(\alpha_1,\dots,\alpha_i-1,\dots,\alpha_n)_g.
\]
\end{thmplaincite}

% [Tutor Note: 原文此处用圆括号 $(\alpha_1,\dots,\alpha_n)_g$ 而非 Witten 惯用的
%  尖括号 $\langle\tau_{\alpha_1}\cdots\tau_{\alpha_n}\rangle_g$。二者只在公式实际使用的
%  顶次情形（括号总次数 $=3g-3+$ 括号内指标个数）相等；对一般多重指标，圆括号是经
%  $C_\alpha$ 归一化后的系数记号。反例：$(g,n)=(1,1)$ 时 $(0)_1=\pi^2/12$，而
%  $\langle\tau_0\rangle_1=0$。本译本照原文保留圆括号（定义见原文 p.31 的
%  $\langle\tau_0^{d_0}\cdots\rangle_g$）。]
%
% [Tutor Note: "dilaton 方程"与"string 方程"是 Witten 猜想（KdV 层级）的两条基本约束，
%  名字来自二维引力的弦论解释：string 方程对应世界面度规的重参数化不变性（$L_{-1}$），
%  dilaton 方程对应伸缩变换（$L_0$）。这两条在原文 p.33 定理 14（[Mir07b, Theorem 6.1]，
%  $L_n(e^F)=0,\ n\ge-1$）中被统一为 $L_{-1}$、$L_0$ 两条约束；p.34 定理 15（[MS08]）
%  给出的是另一套算子 $V_k$ 的版本，勿与此混淆。]
